\IfFileExists{stacks-project.cls}{%
\documentclass{stacks-project}
}{%
\documentclass{amsart}
}

% For dealing with references we use the comment environment
\usepackage{verbatim}
\newenvironment{reference}{\comment}{\endcomment}
%\newenvironment{reference}{}{}
\newenvironment{slogan}{\comment}{\endcomment}
\newenvironment{history}{\comment}{\endcomment}

% For commutative diagrams we use Xy-pic
\usepackage[all]{xy}

% We use 2cell for 2-commutative diagrams.
\xyoption{2cell}
\UseAllTwocells

% We use multicol for the list of chapters between chapters
\usepackage{multicol}

% This is generally recommended for better output
\usepackage{lmodern}
\usepackage[T1]{fontenc}

% For cross-file-references
\usepackage{xr-hyper}

% Package for hypertext links:
\usepackage{hyperref}

% For any local file, say "hello.tex" you want to link to please
% use \externaldocument[hello-]{hello}
\externaldocument[introduction-]{introduction}
\externaldocument[conventions-]{conventions}
\externaldocument[sets-]{sets}
\externaldocument[categories-]{categories}
\externaldocument[topology-]{topology}
\externaldocument[sheaves-]{sheaves}
\externaldocument[sites-]{sites}
\externaldocument[stacks-]{stacks}
\externaldocument[fields-]{fields}
\externaldocument[algebra-]{algebra}
\externaldocument[brauer-]{brauer}
\externaldocument[homology-]{homology}
\externaldocument[derived-]{derived}
\externaldocument[simplicial-]{simplicial}
\externaldocument[more-algebra-]{more-algebra}
\externaldocument[smoothing-]{smoothing}
\externaldocument[modules-]{modules}
\externaldocument[sites-modules-]{sites-modules}
\externaldocument[injectives-]{injectives}
\externaldocument[cohomology-]{cohomology}
\externaldocument[sites-cohomology-]{sites-cohomology}
\externaldocument[dga-]{dga}
\externaldocument[dpa-]{dpa}
\externaldocument[sdga-]{sdga}
\externaldocument[hypercovering-]{hypercovering}
\externaldocument[schemes-]{schemes}
\externaldocument[constructions-]{constructions}
\externaldocument[properties-]{properties}
\externaldocument[morphisms-]{morphisms}
\externaldocument[coherent-]{coherent}
\externaldocument[divisors-]{divisors}
\externaldocument[limits-]{limits}
\externaldocument[varieties-]{varieties}
\externaldocument[topologies-]{topologies}
\externaldocument[descent-]{descent}
\externaldocument[perfect-]{perfect}
\externaldocument[more-morphisms-]{more-morphisms}
\externaldocument[flat-]{flat}
\externaldocument[groupoids-]{groupoids}
\externaldocument[more-groupoids-]{more-groupoids}
\externaldocument[etale-]{etale}
\externaldocument[chow-]{chow}
\externaldocument[intersection-]{intersection}
\externaldocument[pic-]{pic}
\externaldocument[weil-]{weil}
\externaldocument[adequate-]{adequate}
\externaldocument[dualizing-]{dualizing}
\externaldocument[duality-]{duality}
\externaldocument[discriminant-]{discriminant}
\externaldocument[derham-]{derham}
\externaldocument[local-cohomology-]{local-cohomology}
\externaldocument[algebraization-]{algebraization}
\externaldocument[curves-]{curves}
\externaldocument[resolve-]{resolve}
\externaldocument[models-]{models}
\externaldocument[functors-]{functors}
\externaldocument[equiv-]{equiv}
\externaldocument[pione-]{pione}
\externaldocument[etale-cohomology-]{etale-cohomology}
\externaldocument[proetale-]{proetale}
\externaldocument[relative-cycles-]{relative-cycles}
\externaldocument[more-etale-]{more-etale}
\externaldocument[trace-]{trace}
\externaldocument[crystalline-]{crystalline}
\externaldocument[spaces-]{spaces}
\externaldocument[spaces-properties-]{spaces-properties}
\externaldocument[spaces-morphisms-]{spaces-morphisms}
\externaldocument[decent-spaces-]{decent-spaces}
\externaldocument[spaces-cohomology-]{spaces-cohomology}
\externaldocument[spaces-limits-]{spaces-limits}
\externaldocument[spaces-divisors-]{spaces-divisors}
\externaldocument[spaces-over-fields-]{spaces-over-fields}
\externaldocument[spaces-topologies-]{spaces-topologies}
\externaldocument[spaces-descent-]{spaces-descent}
\externaldocument[spaces-perfect-]{spaces-perfect}
\externaldocument[spaces-more-morphisms-]{spaces-more-morphisms}
\externaldocument[spaces-flat-]{spaces-flat}
\externaldocument[spaces-groupoids-]{spaces-groupoids}
\externaldocument[spaces-more-groupoids-]{spaces-more-groupoids}
\externaldocument[bootstrap-]{bootstrap}
\externaldocument[spaces-pushouts-]{spaces-pushouts}
\externaldocument[spaces-chow-]{spaces-chow}
\externaldocument[groupoids-quotients-]{groupoids-quotients}
\externaldocument[spaces-more-cohomology-]{spaces-more-cohomology}
\externaldocument[spaces-simplicial-]{spaces-simplicial}
\externaldocument[spaces-duality-]{spaces-duality}
\externaldocument[formal-spaces-]{formal-spaces}
\externaldocument[restricted-]{restricted}
\externaldocument[spaces-resolve-]{spaces-resolve}
\externaldocument[formal-defos-]{formal-defos}
\externaldocument[defos-]{defos}
\externaldocument[cotangent-]{cotangent}
\externaldocument[examples-defos-]{examples-defos}
\externaldocument[algebraic-]{algebraic}
\externaldocument[examples-stacks-]{examples-stacks}
\externaldocument[stacks-sheaves-]{stacks-sheaves}
\externaldocument[criteria-]{criteria}
\externaldocument[artin-]{artin}
\externaldocument[quot-]{quot}
\externaldocument[stacks-properties-]{stacks-properties}
\externaldocument[stacks-morphisms-]{stacks-morphisms}
\externaldocument[stacks-limits-]{stacks-limits}
\externaldocument[stacks-cohomology-]{stacks-cohomology}
\externaldocument[stacks-perfect-]{stacks-perfect}
\externaldocument[stacks-introduction-]{stacks-introduction}
\externaldocument[stacks-more-morphisms-]{stacks-more-morphisms}
\externaldocument[stacks-geometry-]{stacks-geometry}
\externaldocument[moduli-]{moduli}
\externaldocument[moduli-curves-]{moduli-curves}
\externaldocument[examples-]{examples}
\externaldocument[exercises-]{exercises}
\externaldocument[guide-]{guide}
\externaldocument[desirables-]{desirables}
\externaldocument[coding-]{coding}
\externaldocument[obsolete-]{obsolete}
\externaldocument[fdl-]{fdl}
\externaldocument[index-]{index}

% Theorem environments.
%
\theoremstyle{plain}
\newtheorem{theorem}[subsection]{Theorem}
\newtheorem{proposition}[subsection]{Proposition}
\newtheorem{lemma}[subsection]{Lemma}

\theoremstyle{definition}
\newtheorem{definition}[subsection]{Definition}
\newtheorem{example}[subsection]{Example}
\newtheorem{exercise}[subsection]{Exercise}
\newtheorem{situation}[subsection]{Situation}

\theoremstyle{remark}
\newtheorem{remark}[subsection]{Remark}
\newtheorem{remarks}[subsection]{Remarks}

\numberwithin{equation}{subsection}

% Macros
%
\def\lim{\mathop{\mathrm{lim}}\nolimits}
\def\colim{\mathop{\mathrm{colim}}\nolimits}
\def\Spec{\mathop{\mathrm{Spec}}}
\def\Hom{\mathop{\mathrm{Hom}}\nolimits}
\def\Ext{\mathop{\mathrm{Ext}}\nolimits}
\def\SheafHom{\mathop{\mathcal{H}\!\mathit{om}}\nolimits}
\def\SheafExt{\mathop{\mathcal{E}\!\mathit{xt}}\nolimits}
\def\Sch{\mathit{Sch}}
\def\Mor{\mathop{\mathrm{Mor}}\nolimits}
\def\Ob{\mathop{\mathrm{Ob}}\nolimits}
\def\Sh{\mathop{\mathit{Sh}}\nolimits}
\def\NL{\mathop{N\!L}\nolimits}
\def\CH{\mathop{\mathrm{CH}}\nolimits}
\def\proetale{{pro\text{-}\acute{e}tale}}
\def\etale{{\acute{e}tale}}
\def\QCoh{\mathit{QCoh}}
\def\Ker{\mathop{\mathrm{Ker}}}
\def\Im{\mathop{\mathrm{Im}}}
\def\Coker{\mathop{\mathrm{Coker}}}
\def\Coim{\mathop{\mathrm{Coim}}}

% Boxtimes
%
\DeclareMathSymbol{\boxtimes}{\mathbin}{AMSa}{"02}

%
% Macros for moduli stacks/spaces
%
\def\QCohstack{\mathcal{QC}\!\mathit{oh}}
\def\Cohstack{\mathcal{C}\!\mathit{oh}}
\def\Spacesstack{\mathcal{S}\!\mathit{paces}}
\def\Quotfunctor{\mathrm{Quot}}
\def\Hilbfunctor{\mathrm{Hilb}}
\def\Curvesstack{\mathcal{C}\!\mathit{urves}}
\def\Polarizedstack{\mathcal{P}\!\mathit{olarized}}
\def\Complexesstack{\mathcal{C}\!\mathit{omplexes}}
% \Pic is the operator that assigns to X its picard group, usage \Pic(X)
% \Picardstack_{X/B} denotes the Picard stack of X over B
% \Picardfunctor_{X/B} denotes the Picard functor of X over B
\def\Pic{\mathop{\mathrm{Pic}}\nolimits}
\def\Picardstack{\mathcal{P}\!\mathit{ic}}
\def\Picardfunctor{\mathrm{Pic}}
\def\Deformationcategory{\mathcal{D}\!\mathit{ef}}
\usepackage{cleveref}
\crefname{originalequation}{original equation}{original equations}
\Crefname{originalequation}{Original equation}{Original equations}
\makeatletter
\newlabel{equation-good-diagram}{{0CP8}{source}{}{sourcepointer-1}{}}
\newlabel{equation-good-diagram@cref}{{[originalequation][0][]0CP8}{[1][0][]source}{}{}{}}
\makeatother
\crefname{originallemma}{original lemma}{original lemmas}
\Crefname{originallemma}{Original lemma}{Original lemmas}
\makeatletter
\newlabel{lemma-better-characterize-relative-limit-preserving}{{0CP4}{source}{}{sourcepointer-2}{}}
\newlabel{lemma-better-characterize-relative-limit-preserving@cref}{{[originallemma][0][]0CP4}{[1][0][]source}{}{}{}}
\makeatother
\crefname{originallemma}{original lemma}{original lemmas}
\Crefname{originallemma}{Original lemma}{Original lemmas}
\makeatletter
\newlabel{lemma-descend-isomorphism}{{0852}{source}{}{sourcepointer-3}{}}
\newlabel{lemma-descend-isomorphism@cref}{{[originallemma][0][]0852}{[1][0][]source}{}{}{}}
\makeatother
\crefname{originallemma}{original lemma}{original lemmas}
\Crefname{originallemma}{Original lemma}{Original lemmas}
\makeatletter
\newlabel{lemma-limit-from-good-diagram}{{0CP9}{source}{}{sourcepointer-4}{}}
\newlabel{lemma-limit-from-good-diagram@cref}{{[originallemma][0][]0CP9}{[1][0][]source}{}{}{}}
\makeatother
\crefname{originalproposition}{original proposition}{original propositions}
\Crefname{originalproposition}{Original proposition}{Original propositions}
\makeatletter
\newlabel{proposition-characterize-locally-finite-presentation}{{04AK}{source}{}{sourcepointer-5}{}}
\newlabel{proposition-characterize-locally-finite-presentation@cref}{{[originalproposition][0][]04AK}{[1][0][]source}{}{}{}}
\makeatother
\crefname{originalreference}{original reference}{original references}
\Crefname{originalreference}{Original reference}{Original references}
\makeatletter
\newlabel{situation-limit-noetherian}{{0CP6}{source}{}{sourcepointer-6}{}}
\newlabel{situation-limit-noetherian@cref}{{[originalreference][0][]0CP6}{[1][0][]source}{}{}{}}
\makeatother
\crefname{originalreference}{original reference}{original references}
\Crefname{originalreference}{Original reference}{Original references}
\makeatletter
\newlabel{spaces-morphisms-lemma-diagonal-morphism-finite-type}{{084P}{source}{}{sourcepointer-7}{}}
\newlabel{spaces-morphisms-lemma-diagonal-morphism-finite-type@cref}{{[originalreference][0][]084P}{[1][0][]source}{}{}{}}
\makeatother
\crefname{originalreference}{original reference}{original references}
\Crefname{originalreference}{Original reference}{Original references}
\makeatletter
\newlabel{spaces-morphisms-lemma-factor-factor}{{0CP2}{source}{}{sourcepointer-8}{}}
\newlabel{spaces-morphisms-lemma-factor-factor@cref}{{[originalreference][0][]0CP2}{[1][0][]source}{}{}{}}
\makeatother
\crefname{originalreference}{original reference}{original references}
\Crefname{originalreference}{Original reference}{Original references}
\makeatletter
\newlabel{spaces-morphisms-lemma-properties-diagonal}{{03HK}{source}{}{sourcepointer-9}{}}
\newlabel{spaces-morphisms-lemma-properties-diagonal@cref}{{[originalreference][0][]03HK}{[1][0][]source}{}{}{}}
\makeatother
\title{Morphisms of limits lift to morphisms of systems over good diagrams}
\author{The Stacks Project Authors}
\date{stacks-project:spaces-limits.tex:local-snapshot-20261012; GNU FDL 1.2 or later}
\begin{document}
\maketitle
\noindent\textit{Editorial note.} Counted span is the author's own proof environment, source lines 5359-5401, following the statement at 5323-5357.\par\medskip
\noindent\textit{Editorial note.} No omissions; the cartesian identity for E\_i is displayed as an equation chain.\par\medskip
\noindent\textit{Allowed standard prerequisites (editorial).}\begin{itemize}
\item Limit descriptions from good diagrams. lemma-limit-from-good-diagram (same chapter, cited at \textasciitilde{}5344-5346). X = lim X\_i and Y = lim Y\_i.
\item Characterization of locally finite presentation spaces. proposition-characterize-locally-finite-presentation (same chapter, card L26289, cited at \textasciitilde{}5361-5363). Existence of h : X\_\{i\_0\} -> V.
\item Fibre product criterion for relative limit preservation. lemma-better-characterize-relative-limit-preserving (same chapter, card L26304, cited at \textasciitilde{}5361-5363). The improved form of the proposition applied to V -> B\_\{i\_2\}.
\item Factorization through scheme theoretic images. Morphisms of Spaces, Lemma spaces-morphisms-lemma-factor-factor (cross-chapter, cited at \textasciitilde{}5376-5381). X\_i -> B\_i x\_\{B\_\{i\_2\}\} V factors through Y\_i.
\item Properties of diagonals of finite presentation. Morphisms of Spaces, Lemmas spaces-morphisms-lemma-properties-diagonal and spaces-morphisms-lemma-diagonal-morphism-finite-type (cross-chapter, cited at \textasciitilde{}5388-5391). E\_i -> X\_i is a finitely presented monomorphism.
\item Descent of isomorphisms. lemma-descend-isomorphism (same chapter, cited at \textasciitilde{}5399-5400). E\_i = X\_i for i large, hence f\_i = g\_i.
\end{itemize}\medskip
\ifdefined\hypertarget\hypertarget{sourcepointer-1}{}\fi
\noindent\textit{Source reference.} Original equation 0CP8: sources/stacks-spaces-limits/pinned-original-source.tex, source lines 5264--5264. Navigation pointer only: the counted proof cites this statement by \textbackslash{}ref as a standing input; the deduction the counted proof has to make is carried out inside the counted span, and no part of that deduction is delegated to this pointer.\par\medskip
\ifdefined\hypertarget\hypertarget{sourcepointer-2}{}\fi
\noindent\textit{Source reference.} Original lemma 0CP4: sources/stacks-spaces-limits/pinned-original-source.tex, source lines 1024--1024. Navigation pointer only: the counted proof cites this statement by \textbackslash{}ref as a standing input; the deduction the counted proof has to make is carried out inside the counted span, and no part of that deduction is delegated to this pointer.\par\medskip
\ifdefined\hypertarget\hypertarget{sourcepointer-3}{}\fi
\noindent\textit{Source reference.} Original lemma 0852: sources/stacks-spaces-limits/pinned-original-source.tex, source lines 1925--1925. Navigation pointer only: the counted proof cites this statement by \textbackslash{}ref as a standing input; the deduction the counted proof has to make is carried out inside the counted span, and no part of that deduction is delegated to this pointer.\par\medskip
\ifdefined\hypertarget\hypertarget{sourcepointer-4}{}\fi
\noindent\textit{Source reference.} Original lemma 0CP9: sources/stacks-spaces-limits/pinned-original-source.tex, source lines 5287--5287. Navigation pointer only: the counted proof cites this statement by \textbackslash{}ref as a standing input; the deduction the counted proof has to make is carried out inside the counted span, and no part of that deduction is delegated to this pointer.\par\medskip
\ifdefined\hypertarget\hypertarget{sourcepointer-5}{}\fi
\noindent\textit{Source reference.} Original proposition 04AK: sources/stacks-spaces-limits/pinned-original-source.tex, source lines 694--694. Navigation pointer only: the counted proof cites this statement by \textbackslash{}ref as a standing input; the deduction the counted proof has to make is carried out inside the counted span, and no part of that deduction is delegated to this pointer.\par\medskip
\ifdefined\hypertarget\hypertarget{sourcepointer-6}{}\fi
\noindent\textit{Source reference.} Original reference 0CP6: sources/stacks-spaces-limits/pinned-original-source.tex, source lines 5249--5249. Navigation pointer only: the counted proof cites this statement by \textbackslash{}ref as a standing input; the deduction the counted proof has to make is carried out inside the counted span, and no part of that deduction is delegated to this pointer.\par\medskip
\ifdefined\hypertarget\hypertarget{sourcepointer-7}{}\fi
\noindent\textit{Source reference.} Original reference 084P: focused-seeds/geometry/S9-spaces-limits-morphism-good-diagram/ref-tags.txt, source lines 5--5. Navigation pointer only: the counted proof cites this statement by \textbackslash{}ref as a standing input; the deduction the counted proof has to make is carried out inside the counted span, and no part of that deduction is delegated to this pointer.\par\medskip
\ifdefined\hypertarget\hypertarget{sourcepointer-8}{}\fi
\noindent\textit{Source reference.} Original reference 0CP2: focused-seeds/geometry/S9-spaces-limits-morphism-good-diagram/ref-tags.txt, source lines 6--6. Navigation pointer only: the counted proof cites this statement by \textbackslash{}ref as a standing input; the deduction the counted proof has to make is carried out inside the counted span, and no part of that deduction is delegated to this pointer.\par\medskip
\ifdefined\hypertarget\hypertarget{sourcepointer-9}{}\fi
\noindent\textit{Source reference.} Original reference 03HK: focused-seeds/geometry/S9-spaces-limits-morphism-good-diagram/ref-tags.txt, source lines 7--7. Navigation pointer only: the counted proof cites this statement by \textbackslash{}ref as a standing input; the deduction the counted proof has to make is carried out inside the counted span, and no part of that deduction is delegated to this pointer.\par\medskip
\begin{lemma}
\label{lemma-morphism-good-diagram}
In Situation \ref{situation-limit-noetherian}.
Let $f : X \to Y$ be a morphism of algebraic spaces quasi-separated
and of finite type over $B$. Let
$$
\vcenter{
\xymatrix{
X \ar[r] \ar[d] & W \ar[d] \\
B \ar[r] & B_{i_1}
}
}
\quad\text{and}\quad
\vcenter{
\xymatrix{
Y \ar[r] \ar[d] & V \ar[d] \\
B \ar[r] & B_{i_2}
}
}
$$
be diagrams as in (\ref{equation-good-diagram}). Let
$X = \lim_{i \geq i_1} X_i$ and
$Y = \lim_{i \geq i_2} Y_i$ be the corresponding
limit descriptions as in Lemma \ref{lemma-limit-from-good-diagram}.
Then there exists an $i_0 \geq \max(i_1, i_2)$ and a morphism
$$
(f_i)_{i \geq i_0} : (X_i)_{i \geq i_0} \to (Y_i)_{i \geq i_0}
$$
of inverse systems over $(B_i)_{i \geq i_0}$ such that
such that $f = \lim_{i \geq i_0} f_i$.
If $(g_i)_{i \geq i_0} : (X_i)_{i \geq i_0} \to (Y_i)_{i \geq i_0}$
is a second morphism of inverse systems over $(B_i)_{i \geq i_0}$ such that
such that $f = \lim_{i \geq i_0} g_i$
then $f_i = g_i$ for all $i \gg i_0$.
\end{lemma}

\begin{proof}
Since $V \to B_{i_2}$ is of finite presentation and
$X = \lim_{i \geq i_1} X_i$ we can appeal to Proposition
\ref{proposition-characterize-locally-finite-presentation}
as improved by Lemma \ref{lemma-better-characterize-relative-limit-preserving}
to find an $i_0 \geq \max(i_1, i_2)$ and a morphism $h : X_{i_0} \to V$
over $B_{i_2}$ such that $X \to X_{i_0} \to V$ is equal to $X \to Y \to V$.
For $i \geq i_0$ we get a commutative solid diagram
$$
\xymatrix{
X \ar[d] \ar[r] &
X_i \ar[r] \ar@{..>}[d] \ar@/_2pc/[dd] |!{[d];[ld]}\hole &
X_{i_0} \ar[d]^h \\
Y \ar[r] \ar[d] & Y_i \ar[r] \ar[d] & V \ar[d] \\
B \ar[r] & B_i \ar[r] & B_{i_0}
}
$$
Since $X \to X_i$ has scheme theoretically dense image
and since $Y_i$ is the scheme theoretic image of
$Y \to B_i \times_{B_{i_2}} V$
we find that the morphism $X_i \to B_i \times_{B_{i_2}} V$
induced by the diagram factors through $Y_i$
(Morphisms of Spaces, Lemma \ref{spaces-morphisms-lemma-factor-factor}).
This proves existence.

\medskip\noindent
Uniqueness. Let $E_i \to X_i$ be the equalizer of $f_i$ and $g_i$
for $i \geq i_0$. We have
$E_i = Y_i \times_{\Delta, Y_i \times_{B_i} Y_i, (f_i, g_i)} X_i$.
Hence $E_i \to X_i$ is a monomorphism of finite presentation as a base
change of the diagonal of $Y_i$ over $B_i$, see
Morphisms of Spaces, Lemmas \ref{spaces-morphisms-lemma-properties-diagonal} and
\ref{spaces-morphisms-lemma-diagonal-morphism-finite-type}.
Since $X_i$ is a closed subspace of $B_i \times_{B_{i_0}} X_{i_0}$
and similarly for $Y_i$ we see that
$$
E_i =
X_i \times_{(B_i \times_{B_{i_0}} X_{i_0})} (B_i \times_{B_{i_0}} E_{i_0}) =
X_i \times_{X_{i_0}} E_{i_0}
$$
Similarly, we have $X = X \times_{X_{i_0}} E_{i_0}$. Hence we conclude
that $E_i = X_i$ for $i$ large enough by Lemma \ref{lemma-descend-isomorphism}.
\end{proof}

\end{document}
